% !TeX program = xelatex
% A Working Guide to the Wrong-θ Critique of Conventional Intra-National Distance
% Volume 2: Numerical Evaluation, Empirics, and the Forward Path
% Ian T. Helfrich

\documentclass[11pt,letterpaper,oneside]{book}

% ── Fonts and layout ────────────────────────────────────────────────────────
\usepackage{fontspec}
\usepackage[margin=1in,top=1.05in,bottom=1.05in]{geometry}
\setmainfont{Charter}
\setsansfont{Helvetica Neue}[Scale=0.93]
\setmonofont{Menlo}[Scale=0.82]
\linespread{1.16}
\setlength{\parskip}{0.45em}
\setlength{\parindent}{0pt}

% ── Math and symbols ────────────────────────────────────────────────────────
\usepackage{amsmath,amssymb,amsthm,mathtools,bm}
\usepackage{siunitx}
\usepackage{enumitem}
\setlist{itemsep=0.15em,topsep=0.3em}

\newtheorem{thm}{Theorem}[chapter]
\newtheorem{prop}[thm]{Proposition}
\newtheorem{lem}[thm]{Lemma}
\newtheorem{cor}[thm]{Corollary}
\theoremstyle{definition}
\newtheorem{defn}[thm]{Definition}
\theoremstyle{remark}
\newtheorem*{rem}{Remark}
\newtheorem*{exa}{Example}

\DeclareMathOperator*{\essinf}{ess\,inf}
\DeclareMathOperator*{\esssup}{ess\,sup}
\DeclareMathOperator{\supp}{supp}
\DeclareMathOperator{\vol}{vol}
\newcommand{\R}{\mathbb{R}}
\newcommand{\E}{\mathbb{E}}
\newcommand{\diff}{\mathrm{d}}
\newcommand{\eff}{\mathrm{eff}}
\newcommand{\HM}{\mathrm{HM}}

% ── Colors ──────────────────────────────────────────────────────────────────
\usepackage[dvipsnames,table]{xcolor}
\definecolor{ink}{HTML}{1B1F23}
\definecolor{rust}{HTML}{B5482A}
\definecolor{sage}{HTML}{4C7066}
\definecolor{gold}{HTML}{C28B25}
\definecolor{dim}{HTML}{6C757D}

% ── Boxes ───────────────────────────────────────────────────────────────────
\usepackage[most]{tcolorbox}

\newtcolorbox{insight}{
  enhanced, breakable, sharp corners,
  colback=sage!7, colframe=sage,
  boxrule=0.5pt, left=8pt, right=8pt, top=6pt, bottom=6pt,
  title=Insight, fonttitle=\bfseries\color{sage},
}

\newtcolorbox{notation}{
  enhanced, breakable, sharp corners,
  colback=gold!7, colframe=gold,
  boxrule=0.5pt, left=8pt, right=8pt, top=6pt, bottom=6pt,
  title=Notation parse, fonttitle=\bfseries\color{gold!80!black},
}

\newtcolorbox{worked}{
  enhanced, breakable, sharp corners,
  colback=rust!5, colframe=rust,
  boxrule=0.5pt, left=8pt, right=8pt, top=6pt, bottom=6pt,
  title=Worked example, fonttitle=\bfseries\color{rust},
}

\newtcolorbox{caveat}{
  enhanced, breakable, sharp corners,
  colback=dim!7, colframe=dim,
  boxrule=0.5pt, left=8pt, right=8pt, top=6pt, bottom=6pt,
  title=Caveat, fonttitle=\bfseries\color{dim},
}

% ── Headings ────────────────────────────────────────────────────────────────
\usepackage{titlesec}
\titleformat{\chapter}[display]{\huge\bfseries\color{ink}}{Chapter \thechapter}{0.3em}{\Huge}
\titleformat{\section}{\Large\bfseries\color{ink}}{\thesection}{0.6em}{}
\titleformat{\subsection}{\large\bfseries\color{ink}}{\thesubsection}{0.5em}{}
\titleformat{\part}[display]{\Huge\bfseries\color{rust}}{Part \thepart}{0.5em}{\Huge}

% ── Hyperlinks ──────────────────────────────────────────────────────────────
\usepackage[hidelinks, colorlinks=true,
  linkcolor=rust, citecolor=sage, urlcolor=sage]{hyperref}

% ── ToC ─────────────────────────────────────────────────────────────────────
\usepackage{titletoc}
\setcounter{tocdepth}{1}
\setcounter{secnumdepth}{2}

% ── Listings ────────────────────────────────────────────────────────────────
\usepackage{listings}
\lstdefinestyle{python}{
  basicstyle=\footnotesize\ttfamily,
  language=Python,
  keywordstyle=\color{rust}\bfseries,
  commentstyle=\color{sage}\itshape,
  stringstyle=\color{gold!80!black},
  showstringspaces=false,
  numbers=left, numberstyle=\tiny\color{dim}, numbersep=8pt,
  frame=single, framerule=0.3pt, rulecolor=\color{dim},
  backgroundcolor=\color{white},
  breaklines=true,
  tabsize=4,
}
\lstset{style=python}

% ============================================================================
\begin{document}

% ============================================================================
% TITLE
% ============================================================================
\begin{titlepage}
\centering
\vspace*{2cm}
{\Huge\bfseries A Working Guide to the\\[6pt] Wrong-$\theta$ Critique of\\[4pt] Conventional Intra-National Distance\par}
\vspace{1cm}
{\LARGE\color{rust}\bfseries Volume 2\par}
{\Large Numerical Evaluation, Empirics, and the Forward Path\par}
\vspace{2cm}
{\Large Ian T. Helfrich\par}
\vspace{0.5cm}
{\large Georgia Institute of Technology\par}
\vspace{1.5cm}
{\large \today\par}
\vfill
{\color{dim}\itshape Companion to Volume 1 (Trade Economics Foundations, the Structural CES Gravity Model, Internal Distance, Generalized Means and the Divergence). This volume picks up where Volume 1 ends and develops the numerical, empirical, and methodological apparatus needed to take the wrong-$\theta$ critique into estimation practice.}
\vspace{1cm}
\end{titlepage}

\tableofcontents
\newpage

% ============================================================================
% PREFACE
% ============================================================================
\chapter*{Preface to Volume 2}
\addcontentsline{toc}{chapter}{Preface to Volume 2}

Volume 1 ended on a structural observation: the Head-Mayer (2002) closed-form
expression for intra-national distance has a pole at $\theta = -2$, the
underlying integral diverges in 2D for $\theta \leq -2$, and the
empirically defensible $\theta = \delta(1-\sigma) \in [-7, -3]$ for
$\sigma \in [4, 8]$ falls entirely below the pole. Volume 2 takes that
observation into estimation practice.

The volume has five parts.

\textbf{Part V} (Numerical Evaluation and Regularization) develops the
numerical apparatus. If the integral diverges as a continuous object,
the analyst must choose a finite spatial grain $d_{\min}$ at which to
evaluate the discretized version. The implied $d_{ii}^{\eff}(\theta, d_{\min})$
varies sharply with that choice. Part V documents the regularization
curve on a uniform disk (the cleanest case), then on real country
polygons via the Mollweide projection.

\textbf{Part VI} (Real-Data Construction) walks through the GIS pipeline.
BACI for bilateral trade flows. CEPII Gravity for size, language, contiguity,
and the canonical intra-national distance. GHS-POP, WorldPop, and VIIRS for
population and economic-activity rasters. OpenStreetMap for transport-network
distance. The empirical machinery behind the empirical claim.

\textbf{Part VII} (Regression Specification) is the methodological core.
OLS gravity with origin and destination fixed effects. The perfect-collinearity
trap that corrupts coefficients when log-GDP regressors are included alongside
country FE. The Poisson PML alternative. The Yotov-Piermartini-Monteiro-Larch
pair-year FE convention. A working analyst should leave Part VII with the
ability to estimate a clean gravity regression from raw data.

\textbf{Part VIII} (Empirical Results) walks through the BACI-CEPII
regularization sweep. The Spec A baseline, the Spec B-prime sweep, the
sign-flip threshold by structural $\theta$, the cross-year stability problem.
The interpretation: a coefficient that the literature has treated as a
structural border-effect parameter is in fact a function of an
analyst-chosen discretization grain.

\textbf{Part IX} (Forward Path) sketches what an honest gravity literature
would do. The Anderson-Yotov multilateral-resistance system over time.
Sectoral and pair-year FE specifications. Welfare counterfactuals when
$d_{ii}$ does not exist as a fixed object. The agenda for Paper 5b and beyond.

The volume is designed to be read end-to-end by a working trade economist
who already has Volume 1 internalized. Each chapter has a worked example, a
notation parse on first appearance of a new symbol, and a caveat box where
the assumptions are tighter than they look.

\bigskip
{\color{dim}\itshape A note on style: I am the sole author of this volume.
I write in the first person where my own analytical choices are at stake
(``I take $\delta = 1$ because\ldots'') and in plain declarative voice where
the mathematics is doing the work. The volume is opinionated. When the
literature has been wrong, the volume says so. The reader is welcome to
disagree.}

% ============================================================================
% PART V: NUMERICAL EVALUATION AND REGULARIZATION
% ============================================================================
\part{Numerical evaluation and regularization}

\chapter{Why discretization matters for divergent integrals}
\label{ch:discretization}

\section{The pole, restated}

Volume 1, Chapter 11 established that the 2D moment integral
$\int_0^R r^{\theta+1}\, dr$ diverges at the lower limit for $\theta \leq -2$.
The reason is geometric: the pairwise-distance density of a uniform disk
scales as $f(r) \sim r$ near the origin (this is the 2D version of the
linear vanishing of the pdf at $r = 0$), so the integrand $r^\theta \cdot r$
behaves as $r^{\theta+1}$ near zero, and the integral converges if and
only if $\theta + 1 > -1$, i.e., $\theta > -2$. Below the pole, the
continuous integral has no finite value.

What does this mean operationally? It means that if you write down the
Head-Mayer (2002) closed form
$$d_{ii}(\theta) = \left(\frac{2}{\theta + 2}\right)^{1/\theta} R$$
and try to evaluate it at $\theta = -3$, the formula gives a complex
number on the principal branch. If you write down the integral itself
and try to evaluate it numerically on a discretized region, you get a
finite number, but the number depends on the discretization grain.

This chapter walks through what the discretized number is, how it depends
on the grain, and what the resulting object actually represents.

\section{The discretization}

Suppose we have a uniform-density disk of radius $R$ and we want to
compute the CES expected pairwise distance at $\theta = -5$. The
analytical integral diverges. Numerically, we proceed by sampling
$N$ atoms uniformly within the disk and computing the empirical version:
\begin{equation}
\hat{d}_{ii}^{\eff}(\theta = -5) \;=\; \left( \frac{1}{N(N-1)} \sum_{i \neq j} \|x_i - x_j\|^{-5} \right)^{-1/5}.
\end{equation}
For $N$ atoms in the disk, the smallest pairwise distance is roughly
$R/\sqrt{N}$ (the average nearest-neighbor distance). Atoms closer than
this are rare; the sum is dominated by the few near-collisions.

\begin{notation}
$d_{\min}$ is the regularization cutoff: the spatial grain below which
we either don't sample atoms or we cap pairwise distances. Two natural
implementations: (1) sample $N$ atoms uniformly, accept whatever
$d_{\min}$ falls out; (2) sample many atoms but explicitly cap
$d_{ij} \mapsto \max(d_{ij}, d_{\min})$. Both give the same qualitative
behavior; the second is easier to control.
\end{notation}

\section{What the discretized object measures}

For $\theta > -2$ (above the pole), the empirical estimator
$\hat{d}_{ii}^{\eff}$ converges to the true continuous integral as
$N \to \infty$ or $d_{\min} \to 0$. The discretization grain is a
numerical-accuracy issue, not a conceptual one.

For $\theta \leq -2$, the empirical estimator does not converge. As
$N$ grows (or $d_{\min}$ shrinks), $\hat{d}_{ii}^{\eff}$ decreases
monotonically toward zero. There is no fixed point. The number you
report at any specific $d_{\min}$ is a regularization choice tied to
that grain, not a property of the underlying disk geometry.

\begin{insight}
The continuous CES expected pairwise distance at gravity-consistent
$\theta$ is not a number. It is a function $d_{ii}^{\eff}(\theta, d_{\min})$
defined for every $d_{\min} > 0$, with $d_{ii}^{\eff} \to 0$ as
$d_{\min} \to 0$. The CEPII evaluation at $\theta = -1$ exists because
$\theta = -1$ is above the pole. Below the pole, CEPII would not give a
single number either; it would give a sequence indexed by the city
sample.
\end{insight}

\section{The empirical sweep}

Notebook 19c (the regularization sensitivity script) runs the
discretization at $\theta = -5$ on a uniform disk of radius
$R = 1000$~km, varying $d_{\min} \in \{0.01, 0.1, 1, 5, 10, 50, 100, 200, 500\}$~km.
The output:
\begin{center}
\begin{tabular}{r r r}
\toprule
$d_{\min}$ (km) & $\hat{d}_{ii}^{\eff}$ (km) & ratio to HM at $\theta = -1$ \\
\midrule
0.01 & 0.063 & 0.0001 \\
0.1 & 0.631 & 0.0013 \\
1 & 6.288 & 0.013 \\
5 & 29.4 & 0.059 \\
10 & 51.9 & 0.10 \\
50 & 149.7 & 0.30 \\
100 & 229.4 & 0.46 \\
200 & 352.9 & 0.71 \\
500 & 640.4 & 1.28 \\
\bottomrule
\end{tabular}
\end{center}

The ratio to the HM-at-$\theta = -1$ value ($R/2 = 500$~km) crosses
unity between $d_{\min} = 200$~km and $d_{\min} = 500$~km. At
$d_{\min} = 500$~km, the discretized $\theta = -5$ result exceeds the
$\theta = -1$ analytical result; at $d_{\min} = 0.01$~km, it is four
orders of magnitude smaller. There is no fixed point. The implied
$d_{ii}$ is a function of $d_{\min}$, full stop.

\section{What this implies for empirical practice}

Three operational implications follow.

\textbf{First}, any specific value for $d_{ii}$ reported in a gravity
paper must be accompanied by a statement of the regularization grain.
CEPII's $\theta = -1$ choice produces a single number for each country
because $\theta = -1$ is above the pole. Below the pole, the
single-number convention breaks down.

\textbf{Second}, papers that swap CEPII's distance for a
satellite-calibrated alternative (the Gonchar-Helfrich panel, the
WorldPop-recomputed distances, the OSM-based least-cost-path measures)
are choosing one regularization over another, not solving a measurement
problem. The Paper 5a critique applies to all of them.

\textbf{Third}, the structural-versus-reduced-form distinction in
$\theta$ is not a notational hygiene matter. It is a substantive
methodological choice with empirical consequences. Papers that report
border-puzzle multipliers as structural objects, rather than as
conditional on a regularization, have misrepresented the literature's
empirical content.

\chapter{The regularization sweep on a uniform disk}
\label{ch:disk-sweep}

\section{Setup}

This chapter reproduces the uniform-disk sweep in Python and walks
through the implementation. The script is at
\texttt{notebooks/19c\_regularization\_sensitivity.py} in the project
repository.

The setup: place $N$ atoms uniformly within a disk of radius
$R = 1000$~km, compute all pairwise distances, cap pairwise distances at
$d_{\min}$ from below, then evaluate the CES expected pairwise distance
at $\theta = -5$. Repeat for 5 random seeds, average across seeds, and
sweep $d_{\min}$.

\begin{worked}
The atoms are sampled by inverse-CDF: $r_i = R \sqrt{u_i}$ and
$\phi_i \sim \mathrm{Uniform}(0, 2\pi)$, with $u_i \sim \mathrm{Uniform}(0, 1)$.
The $\sqrt{}$ ensures the resulting density is uniform in 2D area; the
naive sampling $r_i \sim \mathrm{Uniform}(0, R)$ would over-sample
near the center.

Cartesian coordinates: $x_i = r_i \cos \phi_i$, $y_i = r_i \sin \phi_i$.
Pairwise distance: $d_{ij} = \sqrt{(x_i - x_j)^2 + (y_i - y_j)^2}$.

The CES expected distance at $\theta = -5$:
$$\hat{d}_{ii}^{\eff} = \left( \frac{1}{N(N-1)} \sum_{i \neq j} d_{ij}^{-5} \right)^{-1/5}.$$
The capped version: $\tilde{d}_{ij} = \max(d_{ij}, d_{\min})$, then the
same formula.
\end{worked}

\section{Why the result is essentially linear in $d_{\min}$ at low $d_{\min}$}

For $\theta$ well below the pole, the empirical sum is dominated by the
pairs with the smallest distances. With $N$ atoms and a uniform
distribution, the expected number of pairs with $d < d_{\min}$ scales
as $N^2 \cdot (d_{\min}/R)^2$. The sum
$\sum d_{ij}^{-5}$ over these pairs scales as
$N^2 \cdot (d_{\min}/R)^2 \cdot d_{\min}^{-5} = N^2 \cdot d_{\min}^{-3} \cdot R^{-2}$.
Dividing by $N(N-1) \approx N^2$ and taking the $(-1/5)$-th root:
$$\hat{d}_{ii}^{\eff} \sim \left(d_{\min}^{-3} R^{-2}\right)^{-1/5} = d_{\min}^{3/5} R^{2/5}.$$
So $\hat{d}_{ii}^{\eff}$ scales as $d_{\min}^{3/5}$ at low $d_{\min}$:
sub-linear but with a sharp slope. The simulation output (Section 1.4)
fits this scaling.

The same scaling argument generalizes: for $\theta < -2$, the scaling is
$\hat{d}_{ii}^{\eff} \sim d_{\min}^{(\theta+2)/(\theta)}$ at small
$d_{\min}$. At $\theta = -3$: exponent $1/3$. At $\theta = -5$: exponent
$3/5$. At $\theta = -7$: exponent $5/7$. The implied $d_{ii}$ flattens
out (approaches a finite limit only weakly) as $\theta$ becomes more
negative.

\begin{insight}
The slope of $\log \hat{d}_{ii}^{\eff}$ versus $\log d_{\min}$ at small
$d_{\min}$ identifies the scaling exponent
$\alpha(\theta) = (\theta + 2)/\theta$. At $\theta = -1$ (above the pole),
$\alpha = -1$: the implied distance INCREASES as $d_{\min}$ decreases.
At $\theta = -2$ (the pole), $\alpha = 0$: log-flat. At $\theta = -3$,
$\alpha = 1/3$: weak dependence. At $\theta = -7$, $\alpha = 5/7$:
near-linear dependence. The pole at $\theta = -2$ is the sign-change
of the regularization-elasticity.
\end{insight}

\section{What real country polygons add}

The uniform disk is a clean theoretical object. Real countries are not
uniform disks. They have irregular shape (e.g., Norway, Indonesia, Chile),
internal water bodies (e.g., the US Great Lakes), and non-uniform
population density (every country, but extreme cases like Russia and
Canada). The regularization curve on a real country polygon will
generally not match the uniform-disk curve.

Chapter 14 returns to this. For now, the uniform-disk result establishes
the scaling regime; real-country corrections are quantitative deviations
from it, not qualitative departures.

\chapter{Real-country polygons and the Mollweide projection}
\label{ch:mollweide}

\section{Why projections matter for area}

\textbf{[STATUS: skeleton — full development in v0.2 of Volume 2.]}

Pages of text on equal-area projections, the Mollweide ESRI:54009
projection used in the Paper 5a pipeline, comparison with Web Mercator
(common but bad for area), and the resulting area errors for
high-latitude countries. The full chapter would walk through the
\texttt{geopandas.to\_crs("ESRI:54009")} pipeline.

\section{Sampling atoms within a real polygon}

\textbf{[STATUS: skeleton.]}

Rejection sampling: draw points in the polygon's bounding box, reject
those outside the polygon, accept those inside. With $N$ accepted atoms,
the rejection rate equals (bounding box area) / (polygon area). For a
country with low fill factor (e.g., Indonesia with all its archipelagos),
rejection sampling is inefficient. Triangulation-based sampling
(triangulate the polygon, sample uniformly within each triangle weighted
by triangle area) is more efficient.

\section{The country-by-country regularization sweep}

\textbf{[STATUS: skeleton.]}

Full development would tabulate $d_{ii}^{\eff}(\theta = -5, d_{\min})$
for the 27-country pilot subset (USA, RUS, CAN, BRA, AUS, CHN, IND, JPN,
FRA, DEU, NLD, BEL, KOR, GBR, ITA, ESP, MEX, ARG, CHL, EGY, plus 7
smaller countries). The cross-country variation in the regularization
curve tells us how much the small-large country geometry affects the
mechanism.

% ============================================================================
% PART VI: REAL-DATA CONSTRUCTION
% ============================================================================
\part{Real-data construction and the GIS pipeline}

\chapter{BACI and the international trade panel}
\label{ch:baci}

\section{What BACI is and why we use it}

BACI is CEPII's reconciled bilateral trade panel, derived from UN
Comtrade with consistency adjustments. The raw Comtrade data has the
familiar reporting problem: country $i$ records its exports to country
$j$, country $j$ records its imports from country $i$, and the two
numbers disagree (sometimes by factors of 2). BACI applies a
reconciliation algorithm that produces a single best-estimate bilateral
flow $X_{ij}$ for each ordered pair, by year and HS6 product code.

The current versions:
\begin{itemize}
  \item \textbf{BACI HS02 v202401b}: HS 2002 nomenclature, vintage of January
    2024, ``b'' suffix indicating the second revision of that vintage.
    Covers 1995-2022 (most years). What Paper 5a uses.
  \item \textbf{BACI HS17 v202401}: HS 2017 nomenclature, newer products
    (e.g., LNG, electric vehicles, drone components), shorter time series
    (~2018 onward).
  \item \textbf{BACI HS92, HS96, HS07, HS12}: legacy vintages.
\end{itemize}

For a paper that needs a long time series and isn't focused on recent
high-tech goods, HS02 is the right choice. The trade-off is that
HS02's classifications are stale for goods that emerged after 2002.

\section{File structure on disk}

Each year is a separate CSV. The columns (per CEPII's data manual):
\begin{itemize}
  \item \texttt{t}: year (integer)
  \item \texttt{i}: exporter ISO numeric code (3-digit integer; e.g., 840 = USA)
  \item \texttt{j}: importer ISO numeric code
  \item \texttt{k}: HS6 product code (6-digit integer)
  \item \texttt{v}: trade value in current thousand USD
  \item \texttt{q}: quantity (often missing; reported in metric tons or units)
\end{itemize}

A single year's file (e.g., \texttt{BACI\_HS02\_Y2010\_V202401b.csv})
is around 200-300 MB and contains 5-10 million rows. The file is
country-pair-product-year level. For a country-pair-year gravity panel,
you aggregate over the product dimension.

\begin{worked}
\textbf{Building the 2010 country-pair panel from BACI:}

\begin{lstlisting}
import pandas as pd
from pathlib import Path

f = Path("BACI_HS02_Y2010_V202401b.csv")

# Chunked read: aggregate by (exporter, importer) on the fly
pieces = []
for chunk in pd.read_csv(f, usecols=["i", "j", "v"],
                         chunksize=1_000_000):
    # Drop missing trade-value rows
    chunk["v"] = pd.to_numeric(chunk["v"], errors="coerce")
    chunk = chunk.dropna(subset=["v"])
    # Aggregate within chunk
    piece = chunk.groupby(["i", "j"], as_index=False)["v"].sum()
    pieces.append(piece)

# Combine chunks and re-aggregate
agg = pd.concat(pieces).groupby(["i", "j"], as_index=False)["v"].sum()

# Convert to USD (BACI is in 1000 USD)
agg["X_ij_usd"] = agg["v"] * 1000

print(f"{len(agg):,} country-pair flows in 2010")
print(f"Total exports: ${agg['X_ij_usd'].sum() / 1e12:.1f}T")
\end{lstlisting}

For 2010, this gives about 30{,}605 country-pair flows totaling
$\sim\$15$T in world trade.
\end{worked}

\section{Country-code mapping}

BACI uses ISO 3166-1 numeric codes (840 for USA, 156 for China,
etc.). CEPII Gravity uses ISO 3166-1 alpha-3 codes (USA, CHN, etc.).
The mapping is a CEPII-provided lookup table; in practice, you can
build it from the CEPII Gravity file itself:

\begin{lstlisting}
cepii = pd.read_csv("Gravity_V202211.csv")  # via chunked read for size
mapping = cepii[["iso3num_o", "iso3_o"]].dropna().drop_duplicates(
    subset=["iso3num_o"])
num_to_iso = {int(n): s for n, s in mapping.values}

# Apply to BACI
agg["iso3_o"] = agg["i"].map(num_to_iso)
agg["iso3_d"] = agg["j"].map(num_to_iso)
\end{lstlisting}

A small number of BACI rows will have ISO numeric codes that don't
match any CEPII alpha-3 (e.g., disputed territories, historical
codes). Drop these or document them; they amount to less than 1\% of
trade value.

\section{Zero-trade rows}

A key issue for the OLS-versus-PPML choice: the BACI country-pair
panel does NOT include rows where $X_{ij} = 0$. BACI only records
positive trade. To build a panel with zero-trade rows for PPML, you
merge BACI onto the full Cartesian product (CEPII Gravity is already
the full Cartesian product for each year, including zero-trade pairs):

\begin{lstlisting}
# Full CEPII panel for the year
inter = cepii[cepii["iso3_o"] != cepii["iso3_d"]][
    ["iso3_o", "iso3_d", "gdp_o", "gdp_d", "distw_harmonic", "contig"]
]

# Left-join BACI flows; missing => zero trade
inter = inter.merge(agg[["iso3_o", "iso3_d", "X_ij_usd"]],
                    on=["iso3_o", "iso3_d"], how="left")
inter["X_ij_usd"] = inter["X_ij_usd"].fillna(0.0)

print(f"Total rows: {len(inter):,}")
print(f"Zero-trade rows: {(inter['X_ij_usd'] == 0).sum():,} ({100 * (inter['X_ij_usd'] == 0).mean():.1f}%)")
\end{lstlisting}

For 2010, this gives 63{,}234 inter-national rows of which about
32{,}000 (50\%) are zero-trade. The presence of so many zeros is why
PPML is the recommended estimator in the modern gravity literature.

\begin{caveat}
The Wei~(1996) intra-national flow proxy $X_{ii} = \max(0, \mathrm{GDP}_i - \sum_j X_{ij})$
will be zero or close to zero for countries that export more than their
GDP (a sign of re-export or merchanting, e.g., Singapore, Hong Kong,
the Netherlands). These countries get dropped from the intra-national
panel because their Wei proxy is zero. This is a substantive limitation
that the Yotov-Piermartini-Monteiro-Larch (2016) ``trade for production''
proxy partially addresses by computing $X_{ii}$ from sectoral production
rather than total GDP.
\end{caveat}

\section{Aggregating across HS6 products}

For a country-pair-year gravity regression, you typically aggregate
across HS6 products. For a sectoral gravity regression (e.g.,
Caliendo-Parro style), you keep the HS6 dimension and run separate
regressions per sector. BACI is the right source for both.

A nuance: BACI's reconciliation is done at the HS6 level. If you
aggregate to a coarser level (HS2 or HS-section) BEFORE running BACI's
reconciliation, you get different numbers than aggregating AFTER. The
canonical workflow is: keep BACI's HS6-level numbers, aggregate
AFTERWARD if you don't need the sectoral dimension.

\chapter{CEPII Gravity v202211 and the Wei (1996) proxy}
\label{ch:cepii}

\section{What CEPII Gravity contains}

CEPII Gravity is the standard pair-year covariate file for structural
gravity estimation. The current version, v202211 (released November
2022), contains 91 variables over 250 countries from 1948-2020, in a
single Cartesian-product panel of 13.7 million rows (250 origins $\times$
250 destinations $\times$ 60-70 years, with some pair-year gaps).

The variables relevant to Paper 5a:
\begin{itemize}
  \item \texttt{iso3\_o}, \texttt{iso3\_d}: country alpha-3 codes
  \item \texttt{iso3num\_o}, \texttt{iso3num\_d}: country numeric codes
    (for joining to BACI)
  \item \texttt{year}: integer
  \item \texttt{gdp\_o}, \texttt{gdp\_d}: GDP in current USD (origin and
    destination)
  \item \texttt{pop\_o}, \texttt{pop\_d}: population
  \item \texttt{dist}: simple distance between capital cities, km
  \item \texttt{distcap}: distance between capital cities (synonym in some
    vintages)
  \item \texttt{distw\_arithmetic}: arithmetic-mean of city-pairwise
    distances weighted by city population, km. This is the Mayer-Zignago
    (2011) ``distwarith'' variable.
  \item \texttt{distw\_harmonic}: harmonic-mean of the same city-pairwise
    distances. This is the Mayer-Zignago ``distwces'' variable, evaluated
    at $\theta = -1$. \textbf{This is the variable Paper 5a critiques.}
  \item \texttt{contig}: 1 if countries share a land border, 0 otherwise
  \item \texttt{comlang\_off}: 1 if countries share an official language
  \item \texttt{comlang\_ethno}: 1 if at least 9\% of the population speaks
    the same language in both countries
  \item \texttt{col\_dep\_ever}: 1 if there was a colonial relationship
    historically
  \item \texttt{rta}: 1 if there is a regional trade agreement in force in
    that year
\end{itemize}

For the intra-national rows (\texttt{iso3\_o == iso3\_d}), CEPII fills
\texttt{distw\_arithmetic} with the within-country arithmetic-mean of
pairwise city distances and \texttt{distw\_harmonic} with the
within-country harmonic-mean. These are the Head-Mayer (2002) closed
form evaluated at $\theta = +1$ and $\theta = -1$ respectively.

\section{Reading the CSV}

The 13.7M-row file is approximately 4 GB in CSV form. On a FAT32 file
system (e.g., an external USB drive), the file is at or near the
4 GB per-file limit, which can cause read errors. Best practice: store
the file on a local SSD or, even better, convert to Parquet:
\begin{lstlisting}
import pandas as pd
import pyarrow.parquet as pq

# One-time conversion
df = pd.read_csv("Gravity_V202211.csv")
df.to_parquet("Gravity_V202211.parquet", compression="snappy")
\end{lstlisting}

Parquet reads are 5-10x faster and the file is ~700 MB (compressed). For
a year-filtered read (e.g., 2010 only), chunked CSV reads are also
acceptable:
\begin{lstlisting}
def chunked_cepii(year):
    chunks = []
    for chunk in pd.read_csv("Gravity_V202211.csv",
                             dtype={"iso3_o": str, "iso3_d": str, "year": "Int64"},
                             chunksize=500_000, low_memory=False):
        sub = chunk[chunk["year"] == year]
        if len(sub) > 0:
            chunks.append(sub)
    return pd.concat(chunks, ignore_index=True)
\end{lstlisting}

For 2010, a chunked read produces a 63{,}504-row DataFrame (~17 MB) in
about 3-5 minutes on a typical laptop with the CSV on an external drive.

\section{The Wei (1996) intra-national flow proxy}

For intra-national trade (no analog of bilateral exports in trade
statistics), Wei (1996) proposed the residual proxy
\begin{equation}
X_{ii,t} \;=\; \max(0, \, Y_{i,t} - \sum_{j \neq i} X_{ij,t})
\end{equation}
where $Y_{i,t}$ is country $i$'s GDP in year $t$ and the sum is over its
exports to all other countries. The intuition: total GDP equals total
production, which is sold either domestically or as exports, so the
residual is intra-national sales.

The Wei proxy is the standard in the literature but has known limitations.

\begin{caveat}
\textbf{Wei proxy limitations:}
\begin{enumerate}
  \item \textbf{Re-exports and merchanting} (Singapore, Hong Kong,
    Netherlands, Belgium). When a country's exports exceed its
    domestic production (because it re-exports goods produced
    elsewhere), the Wei proxy gives zero or negative. The
    $\max(0, \cdot)$ truncation produces a zero, which excludes the
    country from the intra-national panel.
  \item \textbf{Sectoral composition}. GDP includes services
    (~70\% of US GDP) that are largely non-tradable. The Wei proxy
    treats services as intra-national trade, but they're not really
    ``traded'' in the gravity-model sense.
  \item \textbf{Trade in second-hand goods, intra-firm transfers,
    and missing-flow estimation}. Comtrade doesn't capture all flows.
\end{enumerate}

The Yotov-Piermartini-Monteiro-Larch (2016) ``trade for production''
proxy uses sectoral GDP (specifically, gross output minus exports
within each tradable sector), which addresses the services and
re-export issues. It requires sectoral data that's harder to obtain
but is the modern standard.
\end{caveat}

\section{Building the intra-national rows from CEPII + BACI}

\begin{worked}
\begin{lstlisting}
# Intra-national rows construction
gdp_lookup = dict(
    cepii[["iso3_o", "gdp_o"]]
    .drop_duplicates().dropna().values
)

exports_by_origin = (
    baci.groupby("iso3_o", as_index=False)["X_ij_usd"].sum()
    .set_index("iso3_o")["X_ij_usd"].to_dict()
)

# Pull CEPII's stored intra-national distance
cepii_self = cepii[cepii["iso3_o"] == cepii["iso3_d"]][
    ["iso3_o", "distw_harmonic", "distw_arithmetic"]
].copy()

intra_rows = []
for iso, gdp in gdp_lookup.items():
    if not pd.notna(gdp) or gdp <= 0:
        continue
    exp = exports_by_origin.get(iso, 0.0)
    x_ii = max(0.0, gdp - exp)
    intra_rows.append({
        "iso3_o": iso, "iso3_d": iso,
        "X_ij_usd": x_ii,
        "gdp_o": gdp, "gdp_d": gdp,
        "home": 1,
    })

intra_df = pd.DataFrame(intra_rows)
intra_df = intra_df.merge(cepii_self, on="iso3_o", how="left")
print(f"Intra rows: {len(intra_df):,}")
print(f"With positive Wei proxy: {(intra_df['X_ij_usd'] > 0).sum():,}")
print(f"With CEPII intra distance: {intra_df['distw_harmonic'].notna().sum():,}")
\end{lstlisting}

For 2010, this gives 218 candidate intra rows, of which 10 have BOTH
positive Wei proxy AND non-missing CEPII intra distance. The other
208 are dropped (re-export countries, missing data). The small
$n_{\text{intra}} = 10$ is the source of Paper 5a's identification
caveat in Section 5.3.
\end{worked}

\section{Why the 10 surviving countries are large economies}

The 10 intra-national rows for 2010 typically include: USA, China,
Japan, Germany, France, UK, Italy, Brazil, Canada, Russia. These are
all large economies with GDP well above their export totals. Smaller
countries like the Netherlands (large exports relative to GDP from
re-exports) or open-economy specialists like Singapore drop out.

This selection effect is itself worth thinking about. The
``intra-national'' panel that the border-puzzle literature has been
estimating against is dominated by a few G7 economies. Generalizations
to ``borders matter for trade'' are really ``borders matter for the
particular set of large economies that survive the Wei filter.''

\section{Joining BACI and CEPII into a panel}

The final gravity panel is the union of the inter-national rows
(from CEPII, with $X_{ij}$ filled in from BACI) and the intra-national
rows (constructed from the Wei proxy). For OLS, you keep only positive
$X_{ij}$; for PPML, you keep zero-trade rows.

\begin{lstlisting}
# Inter-national rows
inter = cepii[cepii["iso3_o"] != cepii["iso3_d"]][
    ["iso3_o", "iso3_d", "gdp_o", "gdp_d",
     "distw_harmonic", "contig", "comlang_off"]
].copy()
inter = inter.merge(
    baci[["iso3_o", "iso3_d", "X_ij_usd"]],
    on=["iso3_o", "iso3_d"], how="left"
)
inter["X_ij_usd"] = inter["X_ij_usd"].fillna(0.0)  # zeros for PPML
inter["home"] = 0

# Combine intra + inter
panel = pd.concat([inter, intra_df], ignore_index=True, sort=False)
print(f"Total panel: {len(panel):,} rows")
print(f"Intra-national: {int(panel['home'].sum()):,}")
print(f"Positive flows: {(panel['X_ij_usd'] > 0).sum():,}")
print(f"Zero flows: {(panel['X_ij_usd'] == 0).sum():,}")
\end{lstlisting}

For 2010: 63{,}452 total rows, 10 intra-national, 28{,}053 positive flows
(the OLS sample), 35{,}399 zero flows (PPML uses these too).

\chapter{Population and activity rasters: GHS-POP, WorldPop, VIIRS}
\label{ch:rasters}

\section{Why rasters matter for intra-national distance}

The Mayer-Zignago (2011) CEPII intra-national distance is a
discretization of the underlying CES integral, with the discretization
implemented as a population-weighted harmonic mean over city centroids.
Their city sample comes from a circa-2004 cut of the World Gazetteer
city list, with $\sim 25$ city atoms per country on average and large
coverage gaps. The natural alternative, given modern data, is to
replace the city subsample with a population raster: a regular grid
of cells where each cell carries a population (or activity) weight.
This is the construction at the heart of the Gonchar-Helfrich (2026)
effective-distance panel.

Three rasters are relevant for Paper 5a's empirical strategy.

\section{GHS-POP (JRC Global Human Settlement)}

The \textbf{Global Human Settlement Population grid} is produced by the
European Commission Joint Research Centre. It comes in two
spatial resolutions:
\begin{itemize}
  \item GHS-POP at 100 m: the canonical fine-resolution global
    population raster. Each 100 m $\times$ 100 m cell carries a
    population count.
  \item GHS-POP at 1 km: aggregated version, easier to handle, sufficient
    for country-aggregate analysis.
\end{itemize}

Available years: 1975, 1980, 1985, 1990, 1995, 2000, 2005, 2010, 2015,
2020. The 5-year cadence is fine for cross-section but limiting for
panel work; for year-by-year coverage, WorldPop is the alternative.

The raster is delivered as a GeoTIFF in the Mollweide equal-area
projection (ESRI:54009). The Mollweide choice is important: at
high latitudes, naive lat-lon rasters distort area significantly
(Greenland looks bigger than Africa in Mercator). Mollweide preserves
area exactly, which is what you need when computing population-weighted
spatial statistics.

\begin{worked}
\textbf{Reading GHS-POP and masking by country:}
\begin{lstlisting}
import rasterio
from rasterio.mask import mask
from rasterstats import zonal_stats
import geopandas as gpd

# Country boundaries (GADM, in Mollweide for consistency)
bnd = gpd.read_file("gadm_410_level0.gpkg").to_crs("ESRI:54009")
us = bnd[bnd["GID_0"] == "USA"]

# Mask GHS-POP raster by US polygon
with rasterio.open("GHS_POP_E2020_GLOBE_R2023A_54009_1000_V1_0.tif") as src:
    img, transform = mask(src, us.geometry, crop=True, filled=False)
    nodata = src.nodata or 0
    img = img.filled(nodata)

# Aggregate population
us_pop = img[img > 0].sum()
print(f"US 2020 population: {us_pop / 1e6:.1f}M")  # ~331M
\end{lstlisting}
\end{worked}

\section{WorldPop (annual updates 2000-2024)}

WorldPop is a Southampton-led project providing annual population
rasters at 100 m and 1 km resolution from 2000 onward. The 100 m
version uses a Bayesian dasymetric mapping approach to redistribute
census counts based on building footprints, road density, and
nighttime lights. It is essentially the highest-resolution global
population raster available at annual cadence.

For a Paper 5a-style regularization sweep, WorldPop's annual cadence
is overkill (one year per regression is sufficient). For
Gonchar-Helfrich's 2000-2024 panel, the annual cadence is essential.

\section{VIIRS Day/Night Band (nighttime lights)}

The NOAA VIIRS Day/Night Band on the Suomi-NPP satellite (and later
NOAA-20 and NOAA-21) produces nightly composites of nighttime light
radiance at $\sim 750$ m resolution from 2012 onward. The annual
composites are at \texttt{eogdata.mines.edu}.

Why use VIIRS? Nighttime lights are a well-validated proxy for
economic activity (Henderson, Storeygard, Weil 2012, AER). Where
population rasters tell you where people live, VIIRS tells you where
they're producing, consuming, and lighting up at night. For a
``population-weighted'' versus ``activity-weighted'' comparison of
intra-national distance, you need both.

\section{Working with rasters in Python}

The standard toolchain:
\begin{itemize}
  \item \texttt{rasterio} for reading/writing GeoTIFFs
  \item \texttt{rasterstats.zonal\_stats} for aggregating raster values within polygons
  \item \texttt{geopandas} for vector data (country boundaries, polygons)
  \item \texttt{xarray} + \texttt{rioxarray} for label-aware raster manipulation
\end{itemize}

\begin{caveat}
\textbf{Projection consistency matters.} If your boundaries are in
EPSG:4326 (lat-lon) and your raster is in ESRI:54009 (Mollweide),
\texttt{rasterio.mask.mask} will refuse the operation (or, worse,
silently misalign). Always check \texttt{src.crs} on the raster and
\texttt{gdf.crs} on the vector, and reproject one to match the other.
\end{caveat}

\section{What changes when you use a raster instead of CEPII's cities}

CEPII's city-based intra-national distance, at $\theta = -1$, gives
USA $\approx 1{,}177$ km. The raster-based equivalent, at $\theta = -1$,
on GHS-POP 2020 at 1 km resolution, gives USA $\approx 1{,}140$ km
(a 3\% difference). The two methodologies are operationally similar
\emph{at $\theta = -1$, above the pole}.

\textbf{At $\theta = -5$, the two methodologies diverge sharply.}
The CEPII city sample produces a number (whatever the harmonic-mean
at $\theta = -5$ happens to be on the city subsample). The raster
sample produces a number that depends on the raster resolution.
There is no agreement, because below the pole there is no fixed
truth.

This is the central empirical insight of Paper 5a: the choice of
underlying spatial data (CEPII cities vs. GHS-POP raster vs.\ WorldPop
raster vs.\ a hypothetical even-finer-grain population layer) is the
regularization choice $d_{\min}$. CEPII's cities have effective
$d_{\min} \approx 50-100$ km (city centers are typically separated by
that distance); GHS-POP at 1 km has $d_{\min} = 1$ km; WorldPop at
100 m has $d_{\min} = 0.1$ km. The implied $d_{ii}^{\eff}$ falls
across these choices because the integral has no fixed point.

\chapter{Transport-network distance: OpenStreetMap, ports, and air routes}

\textbf{[STATUS: skeleton.]} Coverage of how to construct a multi-modal
transport graph from OpenStreetMap road data, World Port Index (WPI)
maritime data, and OpenFlights air routes. The Dijkstra least-cost-path
solver in \texttt{networkx} or (faster) Rust via \texttt{paper5-core}.
The Gonchar-Helfrich panel's transport-distance variant.

% ============================================================================
% PART VII: REGRESSION SPECIFICATION
% ============================================================================
\part{Regression specification}

\chapter{OLS gravity with origin and destination fixed effects}
\label{ch:ols-gravity}

\section{The canonical specification}

The workhorse log-OLS gravity regression for cross-section trade data
is:
\begin{equation}
\label{eq:gravity}
\log X_{ij,t} \;=\; \alpha + \beta \log d_{ij,t} + \gamma \cdot \mathbf{1}\{i = j\} + \eta_{i,t} + \mu_{j,t} + \epsilon_{ij,t}.
\end{equation}
The variables:
\begin{itemize}
  \item $X_{ij,t}$: bilateral trade from origin $i$ to destination $j$ in year $t$.
  \item $d_{ij,t}$: bilateral distance (CEPII's \texttt{distw\_harmonic} typically).
  \item $\mathbf{1}\{i=j\}$: the home dummy. 1 for intra-national rows
    (Wei proxy), 0 for inter-national.
  \item $\eta_{i,t}$, $\mu_{j,t}$: origin-year and destination-year FE.
  \item $\epsilon_{ij,t}$: error term, assumed mean-zero conditional on regressors.
\end{itemize}

The coefficient of substantive interest is $\gamma$, the home-bias
multiplier in log units. $e^{\hat{\gamma}}$ is the ratio of
intra-national to equivalent-distance international trade, conditional
on the FE and distance.

\section{What the FE absorb}

In a single-year cross-section, the origin FE $\eta_i$ absorbs:
\begin{itemize}
  \item Origin's GDP (constant per origin in a given year).
  \item Origin's population (constant per origin in a given year).
  \item Origin's average tariff level, average exchange rate, currency
    composition: anything that doesn't vary across destinations within
    the origin.
  \item In a structural reading, the origin's outward multilateral
    resistance term $\Pi_i$.
\end{itemize}

Similarly $\mu_j$ absorbs destination-side characteristics including
the inward multilateral resistance $P_j$. The Anderson-van Wincoop
(2003) insight is that these multilateral resistance terms are why the
McCallum (1995) raw estimate of $22\times$ inflates the structural
border effect; properly absorbed, the multiplier shrinks to $\sim 5\times$.

\begin{insight}
A common student mistake is to include $\log Y_i$ (origin GDP) or
$\log Y_j$ (destination GDP) as explicit regressors alongside the
country FE. In a single-year cross-section, these are perfectly
collinear with the FE structure and should be omitted (or you should
omit the FE and include the GDP terms). Including both is the
collinearity trap of Chapter \ref{ch:collinearity-trap}.
\end{insight}

\section{What the home dummy identifies}

For an intra row (iso\_o = iso\_d = USA, say):
\begin{equation*}
\log \hat{X}_{\text{USA, USA}} = \alpha + \beta \log d_{\text{USA, USA}} + \gamma + \hat{\eta}_{\text{USA}} + \hat{\mu}_{\text{USA}}.
\end{equation*}
For an inter row (iso\_o = USA, iso\_d = CAN):
\begin{equation*}
\log \hat{X}_{\text{USA, CAN}} = \alpha + \beta \log d_{\text{USA, CAN}} + \hat{\eta}_{\text{USA}} + \hat{\mu}_{\text{CAN}}.
\end{equation*}

The home dummy $\gamma$ is identified from the residual difference
between intra-row predictions and inter-row predictions, after
absorbing the origin and destination FE. Specifically, $\hat{\gamma}$
captures the average log-difference between intra-row observed flows
and what the inter-row fit predicts for the intra-row distance.

\section{Why $\hat{\gamma}$ has high variance with small $n_{\text{intra}}$}

In a typical 2010 BACI-CEPII panel, you have $\sim 28{,}000$
inter-national observations and $\sim 10$ intra-national observations.
The FE structure absorbs $\sim 500$ degrees of freedom (250 origins +
250 destinations, minus 1 for the constant). The home dummy is
identified from the intra rows after the FE absorb the country-specific
levels.

With 10 intra observations and 500 FE parameters that can shift to
absorb intra-row residuals, the home dummy has thin identification.
The typical standard error is $\sim 0.7$ in log units, which is
substantial. For Paper 5a, the within-year movement of $\hat{\gamma}$
across $d_{\min}$ is large ($\sim 3$ log units) relative to the
standard error, so the regularization-sensitivity finding is robust.
But the across-year level comparison (year 2010 vs.\ 2015) is not
directly interpretable, because the set of countries that survive
the Wei filter differs across years.

\begin{worked}
\textbf{Estimating the canonical gravity regression in Python:}

\begin{lstlisting}
import pandas as pd
import statsmodels.formula.api as smf
import numpy as np

# panel: pd.DataFrame with columns iso3_o, iso3_d, distance, X_ij_usd, home
df = panel.copy()
df = df[(df["X_ij_usd"] > 0) & (df["distance"] > 0)]
df["log_X"] = np.log(df["X_ij_usd"])
df["log_d"] = np.log(df["distance"])

# Formula API (Patsy handles dummy construction)
formula = "log_X ~ log_d + home + C(iso3_o) + C(iso3_d)"
model = smf.ols(formula, data=df).fit()

print(f"log_d coefficient: {model.params['log_d']:+.4f}")
print(f"home coefficient: {model.params['home']:+.4f}")
print(f"home multiplier: {np.exp(model.params['home']):.2f}x")
print(f"R^2: {model.rsquared:.3f}")
\end{lstlisting}

For year 2010 with CEPII's intra-distance, this gives
$\hat{\beta}_{\log d} = -1.836$, $\hat{\gamma} = +2.132$, multiplier
$8.43\times$. The McCallum-AvW range.
\end{worked}

\section{The case for and against OLS}

OLS works when:
\begin{itemize}
  \item Zero-trade rows are absent (or you're willing to drop them).
  \item The conditional mean is well-described by a log-linear specification.
  \item You're estimating $\E[\log X | Z]$ rather than $\log \E[X | Z]$.
\end{itemize}

OLS fails when:
\begin{itemize}
  \item Zero-trade rows are a substantial fraction (50\% in our 2010
    BACI panel). Dropping them induces a selection bias.
  \item Heteroskedasticity is severe and depends on covariates. The
    Silva-Tenreyro (2006) point: $\Var(X_{ij}|Z)$ scales with
    $E[X_{ij}|Z]^2$ in many gravity settings, which gives biased OLS.
\end{itemize}

PPML (Chapter \ref{ch:ppml}) addresses both failure modes. For Paper 5a's
structural critique, the OLS specification is documented as the
primary; PPML is a robustness check.

\chapter{The perfect-collinearity trap: log-GDP regressors and country FE}

\label{ch:collinearity-trap}

This chapter documents a methodological pitfall I encountered while
running the Paper 5a empirical sweep. It is included in the volume
because the lesson generalizes beyond this specific paper.

\section{The bug}

I built a Spec B-prime regression with the design matrix constructed
manually via \texttt{statsmodels.api.OLS}, including
\texttt{log\_gdp\_o}, \texttt{log\_gdp\_d}, \texttt{contig},
\texttt{comlang\_off} as columns alongside origin and destination
one-hot dummies. The baseline home coefficient came out as $-4.96$ for
year 2010 (multiplier $0.007$, no border puzzle at all), inconsistent
with a known-working version that produced $+2.32$ (multiplier $10$, in
the McCallum-AvW range).

\section{The diagnosis}

\texttt{statsmodels.api.OLS} does NOT automatically drop perfectly
collinear columns. When confronted with a rank-deficient design matrix,
it falls back to the Moore-Penrose pseudo-inverse. The pseudo-inverse
gives a coefficient vector that minimizes the L2-norm subject to the
fit being least-squares optimal, but the coefficient assignment among
collinear columns is arbitrary.

In a single-year cross-section, \texttt{log\_gdp\_o} is perfectly
collinear with the origin one-hot dummy structure (one GDP value per
origin), and \texttt{log\_gdp\_d} with the destination dummies. The
pseudo-inverse redistributes the colinear contribution across the
collinear set, including the home dummy because of numerical noise in
the singular-value decomposition.

\section{The fix}

\texttt{statsmodels.formula.api.ols} uses the formula API (R-style
syntax), which builds the design matrix via Patsy. Patsy detects
perfect collinearity at the dummy-construction stage and drops the
redundant columns. The same regression with the formula API:
\begin{lstlisting}
import statsmodels.formula.api as smf
formula = "log_X ~ log_d + home + C(iso3_o) + C(iso3_d)"
model = smf.ols(formula, data=df).fit()
\end{lstlisting}
gives the correct $+2.13$ home coefficient.

\begin{insight}
\textbf{In single-year cross-sectional regressions with country FE,
do not include log-GDP or other origin-constant / destination-constant
regressors alongside the FE.} The country FE absorb all single-side
variation; the log-GDP terms add no information and corrupt the
non-collinear coefficients via the pseudo-inverse fallback. Use the
formula API or pre-drop the colinear columns. This is not a
\texttt{statsmodels} bug; it is a design choice that requires the user
to handle collinearity explicitly when using \texttt{sm.OLS}.
\end{insight}

\section{Generalizations}

The same trap applies in any cross-sectional regression with high-dim
fixed effects and additional regressors that are perfectly collinear
with the FE structure. Examples:
\begin{itemize}
  \item Country FE plus country-level log-GDP, log-population, capital-output ratio.
  \item Industry FE plus industry-level Herfindahl, average TFP, capital intensity.
  \item Year FE plus year-level interest rate, inflation, recession dummy
    (these vary by year only).
\end{itemize}
The general rule: if a regressor takes one value per FE category, do
not include both. The FE absorbs the regressor.

In panel regressions with multiple years, log-GDP varies within FE
categories (each country has multiple GDP observations across years).
The collinearity is broken. The pitfall is specific to single-year
cross-sections.

\chapter{Poisson PML and the Silva-Tenreyro (2006) zero-trade problem}
\label{ch:ppml}

\section{What Silva-Tenreyro (2006) showed}

The Silva-Tenreyro (2006) ReStat paper is the most-cited methodological
paper in modern gravity estimation. The result, in plain English:
log-OLS gravity is biased when:
\begin{enumerate}
  \item There are zero-trade rows that get dropped from the log
    transformation.
  \item The conditional variance $\Var(X|Z)$ scales with the
    conditional mean $\E[X|Z]$ (typical for trade data).
\end{enumerate}

The bias is upward in absolute value of the distance elasticity:
log-OLS tends to overstate the elasticity of trade with respect to
distance. With zero rows present, the bias can be 30-50\%.

The Silva-Tenreyro fix: estimate
\begin{equation}
\E[X_{ij,t}|Z_{ij,t}] = \exp\left(\alpha + \beta \log d_{ij,t} + \gamma \mathbf{1}\{i=j\} + \eta_{i,t} + \mu_{j,t}\right)
\end{equation}
by Poisson Pseudo-Maximum Likelihood (PPML). The estimator is
consistent under the conditional-mean assumption alone, not requiring
$X_{ij,t}$ to actually follow a Poisson distribution. The
``pseudo-maximum-likelihood'' qualifier captures this: we use the
Poisson likelihood as a quasi-likelihood criterion.

\section{What PPML estimates}

By Jensen's inequality, $\log \E[X|Z] \neq \E[\log X|Z]$ whenever
$\Var(X|Z) > 0$. Log-OLS estimates the right-hand-side conditional
mean; PPML estimates the left-hand-side. These are different objects.
They have different coefficients.

For the gravity context: $\hat{\beta}_{\log d}$ from log-OLS is
typically $-1.5$ to $-2.0$; $\hat{\beta}_{\log d}$ from PPML is
typically $-0.8$ to $-1.0$. The interpretation: trade flows decline
with distance, but the rate of decline depends on what you're
estimating.

\section{The home coefficient under PPML versus OLS}

This is the part that surprised me in Paper 5a's empirical work, so
I'll be explicit. For the same 2010 BACI-CEPII panel:
\begin{itemize}
  \item OLS gives $\hat{\gamma} = +2.13$, multiplier $8.4\times$
    (in the McCallum-AvW range).
  \item PPML (preliminary, with convergence warnings) gives
    $\hat{\gamma} = -5.01$, multiplier $0.007\times$.
\end{itemize}

These are not contradictory; they are estimates of different objects.
PPML estimates the conditional mean of intra-trade given the gravity
covariates; OLS estimates the conditional mean of log-intra-trade.
The Jensen-inequality wedge can be large for high-variance distributions
like trade flows, and the wedge depends on the variance structure.

What survives across both: the regularization-sensitivity pattern. As
$d_{\min}$ shrinks, $\hat{\gamma}$ becomes more negative under both
estimators. The Paper 5a finding ``the home coefficient is regularization-driven''
is robust to estimator choice; the precise magnitude is not.

\section{PPML in \texttt{statsmodels}}

The Python implementation:
\begin{lstlisting}
import statsmodels.formula.api as smf
import statsmodels.api as sm

formula = "X_ij_usd ~ log_d + home + C(iso3_o) + C(iso3_d)"
model = smf.glm(
    formula, data=df,
    family=sm.families.Poisson(link=sm.families.links.Log())
).fit(method="bfgs", maxiter=200, tol=1e-6)
\end{lstlisting}

\begin{caveat}
\textbf{statsmodels.GLM is NOT the right tool for high-dimensional FE
PPML.} With $\sim 500$ FE parameters (one-hot encoded), the BFGS optimizer
struggles to converge. Convergence warnings appear on every cell of a
regularization sweep, and the coefficient estimates are the optimizer's
last iterate rather than the true MLE.

The canonical implementation is the Correia-Guimaraes-Zylkin (2020)
\texttt{ppmlhdfe} Stata package, which uses an absorbing-FE estimator
that re-parameterizes the model in a way that avoids the high-dimensional
design matrix. The Python equivalent is \texttt{pyhdfe} + a custom IRLS
loop, or \texttt{fixest} via \texttt{rpy2}. For Paper 5a, I report
\texttt{statsmodels.GLM} results as preliminary and flag the need for
\texttt{ppmlhdfe} as the natural follow-up.
\end{caveat}

\section{Why this matters for Paper 5a}

Paper 5a's structural claim does not depend on the estimator. The
integration parameter $\theta = \delta(1-\sigma)$ is structural; the
divergence of the integral at $\theta \leq -2$ is mathematical; the
regularization-dependence of any specific number assigned to $d_{ii}$
is a feature of the problem, not of the estimator.

Paper 5a's empirical claim --- the home coefficient depends on the
regularization grain --- is also robust to estimator choice. Both
OLS and PPML produce a monotonic relationship between $\hat{\gamma}$
and $d_{\min}$, with sign and direction matching.

What changes between estimators is the level of $\hat{\gamma}$ at the
CEPII baseline. OLS gives $+2.1$ (in the McCallum range); PPML gives
$-5$ (a sign flip). Both are legitimate estimates of different
objects. The paper reports OLS as primary because it is interpretable
in the McCallum-AvW border-puzzle tradition; PPML is reported as a
robustness check with the appropriate caveats.

\chapter{Pair-year fixed effects: the YPMM cookbook}

\textbf{[STATUS: skeleton.]} The
Yotov-Piermartini-Monteiro-Larch (2016) advanced guide to trade-policy
analysis. The three-way fixed-effect specification:
exporter-year FE, importer-year FE, pair FE. Why pair FE absorbs the
multilateral resistance term. Why this specification is the canonical
2016+ structural-gravity workhorse.

% ============================================================================
% PART VIII: EMPIRICAL RESULTS AND INTERPRETATION
% ============================================================================
\part{The empirical results: Spec B-prime and the regularization sweep}

\chapter{Spec A baseline: replicating the McCallum-AvW finding}

\textbf{[STATUS: skeleton.]} The 2010 panel with CEPII's
\texttt{distw\_harmonic} for both inter and intra. Replicates
McCallum-AvW's $\sim 8\times$ multiplier. The clean baseline against
which all sweeps are compared.

\chapter{Spec B-prime: the regularization sweep results}
\label{ch:bprime}

\section{What Spec B-prime tests}

Spec A (the baseline) uses CEPII's \texttt{distw\_harmonic} for both
intra and inter-national distance. Spec B-prime keeps CEPII's value
for inter-national rows but replaces the intra-national distance with
$d_{ii}^{\eff}(\theta, d_{\min})$ computed numerically on the country
polygon, for $\theta \in \{-3, -5, -7\}$ and
$d_{\min} \in \{0.5, 1, 2, 5, 10, 20, 50, 100, 200, 500\}$~km.

Why hold inter-national distance fixed? The structural critique is
specifically about the intra-national object. Inter-national distance
is also a CES integral in principle, but the bilateral integral has
support disjoint from the divergence at $r = 0$ (Volume 1 Chapter 11),
so it converges at gravity-consistent $\theta$. Swapping only the
intra-national distance isolates the divergence-driven sensitivity.

\section{The full 2010 sweep results}

Three years (2010, 2015, 2020) $\times$ three structural $\theta$'s
$\times$ ten $d_{\min}$ cells $\times$ two baseline specs = 96
regressions. The headline 2010 results:

\begin{center}
\begin{tabular}{r r r r r r r}
\toprule
& \multicolumn{3}{c}{$\hat{\gamma}$} & \multicolumn{3}{c}{$e^{\hat{\gamma}}$} \\
\cmidrule(lr){2-4} \cmidrule(lr){5-7}
$d_{\min}$ (km) & $\theta = -3$ & $\theta = -5$ & $\theta = -7$ & $\theta = -3$ & $\theta = -5$ & $\theta = -7$ \\
\midrule
$0.5$ & $-1.155$ & $-2.919$ & $-3.360$ & $0.32\times$ & $0.05\times$ & $0.03\times$ \\
$1.0$ & $-0.365$ & $-2.054$ & $-2.494$ & $0.69\times$ & $0.13\times$ & $0.08\times$ \\
$2.0$ & $+0.308$ & $-1.293$ & $-1.733$ & $1.36\times$ & $0.27\times$ & $0.18\times$ \\
$5.0$ & $+1.114$ & $-0.378$ & $-0.819$ & $3.05\times$ & $0.68\times$ & $0.44\times$ \\
$10$  & $+1.601$ & $+0.209$ & $-0.231$ & $4.96\times$ & $1.23\times$ & $0.79\times$ \\
$20$  & $+1.976$ & $+0.705$ & $+0.265$ & $7.21\times$ & $2.02\times$ & $1.30\times$ \\
$50$  & $+2.336$ & $+1.282$ & $+0.842$ & $10.3\times$ & $3.60\times$ & $2.32\times$ \\
$100$ & $+2.499$ & $+1.580$ & $+1.140$ & $12.2\times$ & $4.86\times$ & $3.13\times$ \\
$200$ & $+2.593$ & $+1.839$ & $+1.399$ & $13.4\times$ & $6.29\times$ & $4.05\times$ \\
$500$ & $+2.610$ & $+2.031$ & $+1.590$ & $13.6\times$ & $7.62\times$ & $4.91\times$ \\
\bottomrule
\end{tabular}
\end{center}

CEPII baseline at $\theta = -1$ (the conventional choice) gives
$\hat{\gamma} = +2.132$, multiplier $8.43\times$. This sits at the
high-$d_{\min}$ end of the sweep range.

\section{Three observations}

\textbf{First}, the home coefficient moves monotonically with
$d_{\min}$ in every $(\theta, \text{year})$ cell tested. As $d_{\min}$
shrinks (finer spatial resolution), $\hat{\gamma}$ becomes more
negative. As $d_{\min}$ grows (coarser resolution, approaching the
CEPII baseline grain), $\hat{\gamma}$ becomes more positive.

\textbf{Second}, the sign-flip threshold (where $\hat{\gamma}$ crosses
zero) shifts cleanly with $\theta$. At $\theta = -3$, the flip occurs
between $d_{\min} = 1$ and $d_{\min} = 2$~km. At $\theta = -5$, between
$d_{\min} = 5$ and $d_{\min} = 10$~km. At $\theta = -7$, between
$d_{\min} = 10$ and $d_{\min} = 20$~km. More negative $\theta$ means
the formula needs a larger $d_{\min}$ to produce a positive home
coefficient.

\textbf{Third}, the conventional McCallum-Anderson-van~Wincoop border-puzzle
multiplier (the $5\times$ to $22\times$ range that the literature
treats as ``the answer'') sits inside the sweep range and is recovered
at multiple $(\theta, d_{\min})$ combinations. For example, the
$\sim 8\times$ multiplier appears at:
\begin{itemize}
  \item $\theta = -3$, $d_{\min} = 20$~km ($7.21\times$)
  \item $\theta = -5$, $d_{\min} = 500$~km ($7.62\times$)
  \item $\theta = -7$, $d_{\min} > 500$~km (extrapolating)
\end{itemize}
The reduced-form $\theta = -1$ CEPII evaluation produces $8.43\times$
without committing to any structural $\theta$. The literature's
border-puzzle multiplier is not separately identified from the
regularization choice.

\section{Cross-year robustness}

The same sweep on year 2015 ($n_{\text{intra}} = 6$) and year 2020
($n_{\text{intra}} = 8$) produces qualitatively similar patterns:
home moves monotonically more negative as $d_{\min}$ shrinks. The
\emph{levels} differ across years because the small-N intra sample
shifts (different countries pass the Wei filter in different years),
but the within-year pattern is robust.

\begin{insight}
The empirical content of the divergence theorem (Volume 1 Theorem 11.1)
is captured by the slope of $\hat{\gamma}$ versus $\log d_{\min}$. The
slope is steeper for more negative $\theta$, consistent with the
scaling argument in Volume 2 Chapter 2: at $\theta = -3$ the
regularization-elasticity is $1/3$; at $\theta = -7$ it is $5/7$.
The empirical sweep recovers exactly this prediction. The structural
inconsistency in the intra-national distance object propagates into
the home dummy through a mechanism that the closed-form math
predicts.
\end{insight}

\chapter{Cross-year stability and the small-N intra issue}

\textbf{[STATUS: skeleton.]} Why $n_{\text{intra}}$ varies (6 in 2015, 8
in 2020, 10 in 2010). Which countries appear in/out. The implication for
cross-year level comparisons.

\chapter{The qualitative finding: home coefficient as a regularization artifact}
\label{ch:qualitative}

\section{What the empirics show}

Across every $(\theta, \text{year})$ cell in the Paper 5a sweep, the
home-bias coefficient $\hat{\gamma}$ moves monotonically with the
regularization grain $d_{\min}$. At fine resolution (small $d_{\min}$),
$\hat{\gamma}$ is small or negative. At coarse resolution
(large $d_{\min}$, approaching CEPII's effective city-grid grain),
$\hat{\gamma}$ is strongly positive, in the McCallum-Anderson-van~Wincoop
border-puzzle range.

The slope of $\hat{\gamma}$ versus $\log d_{\min}$ is steeper for more
negative $\theta$, consistent with the scaling exponent
$\alpha(\theta) = (\theta + 2)/\theta$ that the divergence theorem
predicts. At gravity-consistent $\theta \in \{-3, -5, -7\}$, the
slope is non-trivial and the sign-flip threshold is in the same
range as CEPII's effective grid resolution.

\section{What this means structurally}

Three structural claims fall out.

\textbf{Claim 1.} The home-bias coefficient in conventional gravity
regressions is not a structural border-effect parameter. It is a
function of the intra-national distance, which is itself a function
of the regularization choice. The literature has been reporting
$\hat{\gamma}$ as if it measured a fixed property of the bilateral
trade environment; it does not.

\textbf{Claim 2.} The McCallum-Anderson-van~Wincoop $5\times$ to
$22\times$ multiplier range corresponds to a specific coarse-resolution
regularization. Finer-resolution regularizations produce smaller
multipliers, and at fine enough resolution, the multiplier is below
$1\times$ (meaning intra-national trade is \emph{less} than what
gravity predicts from inter-national flows at the same distance). The
``border puzzle'' is partly an artifact of the resolution at which
intra-national distance is measured.

\textbf{Claim 3.} Replacing CEPII's static \texttt{distw\_harmonic}
with a satellite-calibrated time-varying alternative (the
Gonchar-Helfrich 2026 panel, WorldPop-based recomputations, OSM-derived
LCP measures) does not solve the underlying problem. It chooses one
regularization over another. The structural inconsistency in the
intra-national distance object propagates regardless of which finite
spatial grain the analyst picks.

\section{What this does NOT mean}

\textbf{This does not mean borders don't matter.} The structural
critique applies to the \emph{magnitude} reported in the literature,
not to the existence of a border effect. A finer-grained empirical
strategy that does not require pinning $d_{ii}$ to a single number
(e.g., the Caliendo-Parro 2015 tetrad estimator, the Anderson-Yotov
multilateral-resistance system) might recover a structural border effect
through a different identification strategy. The Paper 5a critique is
narrow: the specific quantification via CEPII's intra-national distance
is regularization-conditional.

\textbf{This does not mean CEPII's database is wrong.} The harmonic-mean
radial coordinate that CEPII reports at $\theta = -1$ is a perfectly
defensible geometric statistic. What is wrong is the convention of
using that number as if it were the structural CES expected pairwise
distance at the integration parameter implied by empirically defensible
$\sigma$. The two are different objects, and CEPII reports only the
first.

\textbf{This does not mean log-OLS is the wrong estimator.} PPML and
log-OLS give different home coefficients on the same data, but both
exhibit the regularization-sensitivity pattern. The estimator choice
is independent of the structural critique.

\section{What an honest gravity literature would do}

Three operational recommendations.

\textbf{First}, report intra-national distance with its spatial
resolution explicitly. A paper that uses CEPII's value should
acknowledge that the value corresponds to a $\sim 50-100$ km effective
grid. A paper that uses WorldPop should acknowledge the 100 m or 1 km
resolution. The reported home-bias multiplier is conditional on this
choice.

\textbf{Second}, report a sensitivity analysis. The Paper 5a sweep
table is a template: $\hat{\gamma}$ across $d_{\min}$ at the
structurally consistent $\theta$. The reader can then see the
regularization dependence and judge the structural content of the
reported multiplier.

\textbf{Third}, move toward identification strategies that do not
require a single number for $d_{ii}$. The Caliendo-Parro tetrad
estimator. The Eaton-Kortum $d_{ii} = 1$ assumption (which is
structurally consistent: ``intra-national distance equals 1'' just
means trade costs are normalized to 1 at home). The
Costinot-Donaldson-Komunjer Ricardian framework with importer-exporter
FE absorbing all bilateral trade costs. These are not workarounds;
they are honest acknowledgments that the structural integral does not
exist as a fixed object when the structural $\theta$ is in its
empirically defensible range.

\section{The agenda for Paper 5b}

The follow-up paper (Paper 5b, in concept) would develop one of these
identification strategies in detail. The most promising candidates:

\textbf{Option A: Anderson-Yotov MR with explicit regularization-grain
documentation.} Estimate the structural gravity model with the
Anderson-Yotov multilateral-resistance system (which uses no separate
intra-national distance), and document how the MR-implied border
effect depends on the regularization grain of the bilateral distance
inputs. This is the most direct extension of Paper 5a.

\textbf{Option B: Caliendo-Parro tetrad on BACI sectoral data.} The
tetrad estimator cancels distance entirely; the structural-vs-reduced-form
$\theta$ conflation does not arise. The paper would estimate sectoral
tetrads on BACI HS6 data and quantify the border effect via a different
identification strategy.

\textbf{Option C: A regularization-aware structural estimator.}
Develop an estimator that takes $d_{\min}$ as an input parameter and
returns conditional structural estimates. The analyst would then report
results as a function of $d_{\min}$, not as a single number.

Paper 5b is concept-stage. The technical details depend on which of
these options proves most tractable. Paper 5a's contribution is the
diagnosis; Paper 5b's would be the constructive proposal.

% ============================================================================
% PART IX: FORWARD PATH
% ============================================================================
\part{Forward path}

\chapter{Anderson-Yotov multilateral resistance over time}

\textbf{[STATUS: skeleton.]} The 2010 Anderson-Yotov paper on the
multilateral-resistance term as a time-varying object. How the
regularization issue interacts with the MR system.

\chapter{Sectoral and pair-year FE specifications}

\textbf{[STATUS: skeleton.]} What changes when you move to sectoral
trade data. Whether the regularization-sensitivity result survives the
YPMM three-way FE specification.

\chapter{Welfare counterfactuals when $d_{ii}$ does not exist}

\textbf{[STATUS: skeleton.]} The Caliendo-Parro (2015) exact-hat
methodology. Welfare effects of trade-policy changes computed without
a single number for $d_{ii}$. The forward path for Paper 5b.

\chapter{What an honest gravity literature would do}

\textbf{[STATUS: skeleton.]} The constructive proposal. Report
border-puzzle multipliers as conditional on a regularization. Acknowledge
the structural-versus-reduced-form $\theta$ distinction. Move to
estimation strategies that do not require pinning $d_{ii}$ to a single
number.

% ============================================================================
% BACK MATTER
% ============================================================================
\appendix
\chapter{Notation key (Volume 2 additions)}

\begin{small}
\begin{tabular}{l l l}
\toprule
Symbol & Meaning & Units \\
\midrule
$N$ & Number of atoms in discretization & dimensionless \\
$d_{\min}$ & Regularization cutoff & km \\
$\hat{d}_{ii}^{\eff}$ & Discretized CES expected pairwise distance & km \\
$\alpha(\theta)$ & Regularization scaling exponent & dimensionless \\
$\mathrm{ESRI:54009}$ & Mollweide equal-area projection & --- \\
$X_{ij,t}$ & Bilateral trade flow in year $t$ & USD \\
$\eta_{i,t}$ & Origin-year fixed effect & log units \\
$\mu_{j,t}$ & Destination-year fixed effect & log units \\
\bottomrule
\end{tabular}
\end{small}

\chapter{Code listings index}

The notebooks and modules referenced in Volume 2:
\begin{itemize}
  \item \texttt{notebooks/19c\_regularization\_sensitivity.py} — Chapter \ref{ch:disk-sweep}
  \item \texttt{notebooks/18d\_spec\_b\_prime\_fixed.py} — Chapter \ref{ch:collinearity-trap}
  \item \texttt{notebooks/18e\_spec\_b\_prime\_ppml.py} — PPML robustness check
  \item \texttt{src/paper5/region\_raster.py} — Country polygon to RasterDist
  \item \texttt{src/paper5/ot\_distance.py} — CES generalized-mean distance
\end{itemize}

\end{document}
